\chapter{高等数学：多元微分学}

\section{复合函数求导与全微分}

\begin{problemblock}
\textbf{例1.} 设
\[
u(x,y)=\varphi(x+y)+\varphi(x-y)+\int_{x-y}^{x+y}\psi(t)\,dt,
\]
其中函数 \(\varphi\) 具有二阶导数，\(\psi\) 具有一阶导数，则必有（\quad）

A. \(\displaystyle \frac{\partial^2u}{\partial x^2}=-\frac{\partial^2u}{\partial y^2}\)
\quad
B. \(\displaystyle \frac{\partial^2u}{\partial x^2}=\frac{\partial^2u}{\partial y^2}\)

C. \(\displaystyle \frac{\partial^2u}{\partial x\partial y}=\frac{\partial^2u}{\partial y^2}\)
\quad
D. \(\displaystyle \frac{\partial^2u}{\partial x\partial y}=\frac{\partial^2u}{\partial x^2}\)
\end{problemblock}
\begin{solutionblock}
\analysis{令 \(s=x+y,\ r=x-y\)。题中每一项都是 \(s\) 的函数加 \(r\) 的函数，因此满足一维波动方程 \(u_{xx}=u_{yy}\)。}
设
\[
F'(t)=\psi(t).
\]
则
\[
\int_{x-y}^{x+y}\psi(t)\,dt=F(x+y)-F(x-y).
\]
所以
\[
u(x,y)=\bigl[\varphi(s)+F(s)\bigr]+\bigl[\varphi(r)-F(r)\bigr],
\quad s=x+y,\ r=x-y.
\]
对形如 \(A(s)+B(r)\) 的函数，
\[
u_{xx}=A''(s)+B''(r),\qquad
u_{yy}=A''(s)+B''(r).
\]
因此
\[
u_{xx}=u_{yy}.
\]
\answer{B}
\examnote{看到 \(x+y,x-y\) 的组合，优先令 \(s=x+y,r=x-y\)，常会出现 \(u_{xx}=u_{yy}\)。}
\end{solutionblock}

\begin{problemblock}
\textbf{例2.} 设 \(z=f(x,y)\) 由
\[
f(x+y,x-y)=x^2-y^2-xy
\]
确定，求 \(dz\)。
\end{problemblock}
\begin{solutionblock}
\analysis{先把 \(f(X,Y)\) 显式写出，再求全微分。}
令
\[
X=x+y,\qquad Y=x-y.
\]
则
\[
x=\frac{X+Y}{2},\qquad y=\frac{X-Y}{2}.
\]
于是
\[
x^2-y^2=XY,\qquad
xy=\frac{X^2-Y^2}{4}.
\]
故
\[
f(X,Y)=XY-\frac{X^2-Y^2}{4}
=XY-\frac14X^2+\frac14Y^2.
\]
把 \(X,Y\) 改回 \(x,y\)，有
\[
z=f(x,y)=xy-\frac14x^2+\frac14y^2.
\]
因此
\[
dz=\left(y-\frac x2\right)dx+\left(x+\frac y2\right)dy.
\]
\answer{\(\displaystyle dz=\left(y-\frac x2\right)dx+\left(x+\frac y2\right)dy\)}
\examnote{给出 \(f(x+y,x-y)\) 时，不要直接对原式乱套链式法则；先做变量反解最稳。}
\end{solutionblock}

\begin{problemblock}
\textbf{例3.} 设函数 \(f(x,y)\) 可微，且
\[
f(x+1,e^x)=x(x+1)^2,\qquad
f(x,x^2)=2x^2\ln x,
\]
则 \(df(1,1)=(\quad)\)

A. \(dx+dy\)\qquad B. \(dx-dy\)\qquad C. \(dy\)\qquad D. \(-dy\)
\end{problemblock}
\begin{solutionblock}
\analysis{两条曲线都经过 \((1,1)\)，分别沿曲线求导，得到 \(f_x(1,1),f_y(1,1)\) 的两个线性方程。}
对
\[
f(x+1,e^x)=x(x+1)^2
\]
在 \(x=0\) 处求导，得
\[
f_x(1,1)+f_y(1,1)=1.
\]
对
\[
f(x,x^2)=2x^2\ln x
\]
在 \(x=1\) 处求导，得
\[
f_x(1,1)+2f_y(1,1)=2.
\]
解得
\[
f_x(1,1)=0,\qquad f_y(1,1)=1.
\]
所以
\[
df(1,1)=dy.
\]
\answer{C}
\examnote{已知复合函数在曲线上的表达式，就是给方向导数信息；交点处列方程即可。}
\end{solutionblock}

\begin{problemblock}
\textbf{例4.} 设函数
\[
z=f\left(2x-y,\frac yx\right)
\]
具有连续偏导数，且
\[
dz\big|_{(1,3)}=-7\,dx+3\,dy,
\]
则
\[
f'_1(-1,3)+f'_2(-1,3)=\underline{\hspace{2cm}}.
\]
\end{problemblock}
\begin{solutionblock}
\analysis{在 \((x,y)=(1,3)\) 时，内层变量为 \((-1,3)\)。用链式法则列出 \(z_x,z_y\)。}
记
\[
A=f'_1(-1,3),\qquad B=f'_2(-1,3).
\]
因
\[
u=2x-y,\qquad v=\frac yx,
\]
在 \((1,3)\) 处
\[
u_x=2,\quad v_x=-\frac y{x^2}=-3,\qquad
u_y=-1,\quad v_y=\frac1x=1.
\]
由 \(dz=-7dx+3dy\)，得
\[
z_x=2A-3B=-7,
\]
\[
z_y=-A+B=3.
\]
解得
\[
A=-2,\qquad B=1.
\]
因此
\[
f'_1(-1,3)+f'_2(-1,3)=-1.
\]
\answer{\(-1\)}
\examnote{复合求导题要先写清内层变量在给定点的取值，再把偏导数记成未知量。}
\end{solutionblock}

\begin{problemblock}
\textbf{例5.} 设函数 \(f(u,v)\) 可微，且
\[
\frac{\partial f(u,v)}{\partial u}-2\frac{\partial f(u,v)}{\partial v}
=\cos(u+v).
\]
设
\[
g(x,y)=f(x,y-2x).
\]
\begin{enumerate}[label=(\arabic*), leftmargin=2.2em]
\item 求 \(\dfrac{\partial g(x,y)}{\partial x}\)；
\item 若 \(g(0,y)=y^2-\sin y\)，求 \(g(x,y)\) 与 \(f(u,v)\) 的表达式。
\end{enumerate}
\end{problemblock}
\begin{solutionblock}
\analysis{由于 \(g_x=f_u-2f_v\)，题设的偏微分方程正好直接给出 \(g_x\)。}
(1) 由链式法则，
\[
g_x(x,y)=f_u(x,y-2x)-2f_v(x,y-2x).
\]
题设给出
\[
f_u(u,v)-2f_v(u,v)=\cos(u+v).
\]
取 \(u=x,\ v=y-2x\)，得
\[
g_x(x,y)=\cos(y-x).
\]

(2) 对 \(x\) 积分：
\[
g(x,y)=-\sin(y-x)+C(y).
\]
由
\[
g(0,y)=y^2-\sin y
\]
得 \(C(y)=y^2\)。所以
\[
g(x,y)=y^2-\sin(y-x).
\]
又
\[
g(x,y)=f(x,y-2x).
\]
令 \(u=x,\ v=y-2x\)，则 \(y=v+2u\)，故
\[
f(u,v)=(v+2u)^2-\sin(u+v).
\]
\answer{(1) \(\cos(y-x)\)；(2) \(g(x,y)=y^2-\sin(y-x)\)，\(f(u,v)=(v+2u)^2-\sin(u+v)\)。}
\examnote{这类题的设计点通常是让链式法则中的组合恰好等于题设微分算子。}
\end{solutionblock}

\begin{problemblock}
\textbf{例6.} 设 \(y=g(x,z)\) 与 \(z=z(x,y)\) 是由方程
\[
f(x-z,xy)=0
\]
确定的函数，求 \(\dfrac{dy}{dx}\)。
\end{problemblock}
\begin{solutionblock}
\analysis{把 \(z\) 看作参数，\(y\) 是 \(x\) 的隐函数，对 \(f(x-z,xy)=0\) 关于 \(x\) 求导。}
记
\[
u=x-z,\qquad v=xy.
\]
当 \(z\) 固定、\(y=y(x)\) 时，
\[
\frac{du}{dx}=1,\qquad
\frac{dv}{dx}=y+x\frac{dy}{dx}.
\]
对
\[
f(u,v)=0
\]
求导得
\[
f_1'(u,v)\cdot1+f_2'(u,v)\left(y+x\frac{dy}{dx}\right)=0.
\]
因此
\[
\frac{dy}{dx}
=-\frac{f_1'(x-z,xy)+y f_2'(x-z,xy)}{x f_2'(x-z,xy)}.
\]
\answer{\(\displaystyle \frac{dy}{dx}
=-\frac{f_1'(x-z,xy)+y f_2'(x-z,xy)}{x f_2'(x-z,xy)}\)}
\examnote{隐函数求导时，先明确谁是自变量、谁是因变量、谁暂作参数。}
\end{solutionblock}

\begin{problemblock}
\textbf{例7.} 已知函数
\[
f(x,y,z)=x^3y^2z
\]
及方程
\[
x+y+z-3+e^{-3}=e^{-(x+y+z)}\tag{*}
\]
\begin{enumerate}[label=(\arabic*), leftmargin=2.2em]
\item 如果 \(x=x(y,z)\) 是由方程 \((*)\) 确定的隐函数，满足 \(x(1,1)=1\)，又
\[
u=f(x(y,z),y,z),
\]
求 \(\left.\dfrac{\partial u}{\partial y}\right|_{(1,1,1)}\)；
\item 如果 \(z=z(x,y)\) 是由方程 \((*)\) 确定的隐函数，满足 \(z(1,1)=1\)，又
\[
w=f(x,y,z(x,y)),
\]
求 \(\left.\dfrac{\partial w}{\partial y}\right|_{(1,1,1)}\)。
\end{enumerate}
\end{problemblock}
\begin{solutionblock}
\analysis{方程只含 \(s=x+y+z\)。在 \((1,1,1)\) 附近由单调性可知 \(s=3\)，所以隐函数的相关偏导都是 \(-1\)。}
令
\[
H(s)=s-3+e^{-3}-e^{-s}.
\]
方程 \((*)\) 为 \(H(x+y+z)=0\)。由于
\[
H'(s)=1+e^{-s}>0,
\]
在 \(s=3\) 附近只能有
\[
x+y+z=3.
\]
(1) 若 \(x=x(y,z)\)，则
\[
x_y=-1.
\]
又
\[
f_x=3x^2y^2z,\qquad f_y=2x^3yz.
\]
在 \((1,1,1)\) 处，
\[
f_x=3,\qquad f_y=2.
\]
所以
\[
u_y=f_xx_y+f_y=3(-1)+2=-1.
\]

(2) 若 \(z=z(x,y)\)，则
\[
z_y=-1.
\]
又
\[
f_y=2x^3yz,\qquad f_z=x^3y^2.
\]
在 \((1,1,1)\) 处，
\[
f_y=2,\qquad f_z=1.
\]
所以
\[
w_y=f_y+f_zz_y=2+1\cdot(-1)=1.
\]
\answer{(1) \(-1\)；(2) \(1\)。}
\examnote{若隐方程只依赖 \(x+y+z\)，先化成和为常数，可大幅减少计算。}
\end{solutionblock}

\section{隐函数、多元极限与偏导}

\begin{problemblock}
\textbf{例8.} 设函数 \(f(x,y)\) 在 \((0,1)\) 的某领域内一阶偏导数连续，
\[
f(0,1)=0,\qquad f_y'(0,1)=1.
\]
则由
\[
f\left(x,\int_1^t\ln s\,ds\right)=0
\]
可得（\quad）

A. 在点 \((0,1)\) 附近可确定 \(t=t(x)\)，且 \(t'(0)=-f_x'(0,1)\)

B. 在点 \((0,1)\) 附近可确定 \(t=t(x)\)，且 \(t'(0)=-1\)

C. 在点 \((0,e)\) 附近可确定 \(t=t(x)\)，且 \(t'(0)=-f_x'(0,1)\)

D. 在点 \((0,e)\) 附近可确定 \(t=t(x)\)，且 \(t'(0)=-1\)
\end{problemblock}
\begin{solutionblock}
\analysis{要让 \(f\) 的第二个变量等于 \(1\)，需找 \(\int_1^t\ln s\,ds=1\)，这给出 \(t=e\)。}
计算
\[
\int_1^t\ln s\,ds=t\ln t-t+1.
\]
当 \(t=e\) 时，该积分为
\[
e\cdot1-e+1=1.
\]
所以对应点为 \((x,t)=(0,e)\)，不是 \((0,1)\)。

令
\[
F(x,t)=f\left(x,\int_1^t\ln s\,ds\right).
\]
则
\[
F(0,e)=f(0,1)=0.
\]
并且
\[
F_t(0,e)=f_y'(0,1)\ln e=1\ne0.
\]
由隐函数定理，\(t\) 可在 \((0,e)\) 附近表示为 \(t=t(x)\)。又
\[
t'(0)=-\frac{F_x(0,e)}{F_t(0,e)}
=-\frac{f_x'(0,1)}{1}
=-f_x'(0,1).
\]
\answer{C}
\examnote{隐函数题先找准基点；本题的基点由积分上限 \(t=e\) 决定。}
\end{solutionblock}

\begin{problemblock}
\textbf{例9.} 计算
\[
\lim_{(x,y)\to(0,0)}
\left(2^{x^2+y^2}+xy^2\right)^{\frac{x}{x^2+y^2}}.
\]
\end{problemblock}
\begin{solutionblock}
\analysis{幂指函数先取对数。记 \(r^2=x^2+y^2\)，看指数乘以底数对数的极限。}
设
\[
L=\left(2^{x^2+y^2}+xy^2\right)^{\frac{x}{x^2+y^2}},
\qquad r^2=x^2+y^2.
\]
取对数：
\[
\ln L=\frac{x}{r^2}\ln\left(2^{r^2}+xy^2\right).
\]
当 \(r\to0\) 时，
\[
2^{r^2}=1+(\ln2)r^2+O(r^4),
\]
且
\[
xy^2=O(r^3).
\]
因此
\[
\ln\left(2^{r^2}+xy^2\right)
=(\ln2)r^2+O(r^3).
\]
于是
\[
\ln L
=\frac{x}{r^2}\bigl[(\ln2)r^2+O(r^3)\bigr]
=x\ln2+O(xr)\to0.
\]
故
\[
L\to e^0=1.
\]
\answer{\(1\)}
\examnote{多元幂指极限常用“取对数 + 估阶”，尤其要注意 \(xy^2=O(r^3)\)。}
\end{solutionblock}

\begin{problemblock}
\textbf{例10.} 设
\[
f(x,y)=\frac{x^2-y^2}{x^2+y^2},
\]
\[
I_1=\lim_{x\to0}\left[\lim_{y\to0}f(x,y)\right],\quad
I_2=\lim_{y\to0}\left[\lim_{x\to0}f(x,y)\right],\quad
I_3=\lim_{\substack{x\to0\\y\to0}}f(x,y).
\]
则（\quad）

A. \(I_1,I_2\) 存在，\(I_3\) 不存在
\quad
B. \(I_1,I_2,I_3\) 均不存在

C. \(I_1,I_2,I_3\) 均存在
\quad
D. \(I_1,I_2\) 不存在，\(I_3\) 存在
\end{problemblock}
\begin{solutionblock}
\analysis{累次极限与二重极限不同。先固定一个变量，再令另一个变量趋于零。}
对 \(I_1\)，先令 \(y\to0\)。当 \(x\ne0\) 固定时，
\[
\lim_{y\to0}\frac{x^2-y^2}{x^2+y^2}=1.
\]
再令 \(x\to0\)，得
\[
I_1=1.
\]
对 \(I_2\)，先令 \(x\to0\)。当 \(y\ne0\) 固定时，
\[
\lim_{x\to0}\frac{x^2-y^2}{x^2+y^2}=-1.
\]
再令 \(y\to0\)，得
\[
I_2=-1.
\]
但二重极限不存在，因为沿 \(y=0\) 趋近得 \(1\)，沿 \(x=0\) 趋近得 \(-1\)。
\answer{A}
\examnote{累次极限存在不保证二重极限存在；路径检验仍是二重极限的核心。}
\end{solutionblock}

\begin{problemblock}
\textbf{例11.} 已知
\[
f(x,y)=
\begin{cases}
(x^2+y^2)\sin\dfrac1{xy},&xy\ne0,\\
0,&xy=0.
\end{cases}
\]
求 \(\dfrac{\partial f}{\partial x}\)。
\end{problemblock}
\begin{solutionblock}
\analysis{在 \(xy\ne0\) 处直接求导；在坐标轴上必须回到偏导定义。}
当 \(xy\ne0\) 时，
\[
f_x
=2x\sin\frac1{xy}
+(x^2+y^2)\cos\frac1{xy}\left(-\frac1{x^2y}\right),
\]
即
\[
f_x=2x\sin\frac1{xy}
-\frac{x^2+y^2}{x^2y}\cos\frac1{xy}.
\]
当 \(y=0\) 时，
\[
f_x(x,0)=\lim_{h\to0}\frac{f(x+h,0)-f(x,0)}h=0.
\]
特别地 \(f_x(0,0)=0\)。

当 \(x=0,\ y\ne0\) 时，
\[
\frac{f(h,y)-f(0,y)}h
=\frac{h^2+y^2}{h}\sin\frac1{hy},
\]
该极限不存在，所以 \(f_x(0,y)\) 不存在。
\answer{\[
f_x(x,y)=
\begin{cases}
2x\sin\dfrac1{xy}
-\dfrac{x^2+y^2}{x^2y}\cos\dfrac1{xy},&xy\ne0,\\
0,&y=0,\\
\text{不存在},&x=0,\ y\ne0.
\end{cases}
\]}
\examnote{分段多元函数在坐标轴上的偏导不能套非轴区域公式，必须用定义。}
\end{solutionblock}

\begin{problemblock}
\textbf{例12.} 已知
\[
f(x,y)=
\begin{cases}
(x^2+y^2)\sin\dfrac1{xy},&xy\ne0,\\
0,&xy=0,
\end{cases}
\]
则在点 \((0,0)\) 处（\quad）

A. \(\dfrac{\partial f}{\partial x}\) 连续，\(f(x,y)\) 可微

B. \(\dfrac{\partial f}{\partial x}\) 连续，\(f(x,y)\) 不可微

C. \(\dfrac{\partial f}{\partial x}\) 不连续，\(f(x,y)\) 可微

D. \(\dfrac{\partial f}{\partial x}\) 不连续，\(f(x,y)\) 不可微
\end{problemblock}
\begin{solutionblock}
\analysis{可微看 \(f(x,y)=o(\sqrt{x^2+y^2})\)；偏导连续看例11中的 \(f_x\) 是否趋于 \(0\)。}
先看可微性。因为
\[
|f(x,y)|\le x^2+y^2,
\]
故
\[
\frac{|f(x,y)-f(0,0)|}{\sqrt{x^2+y^2}}
\le\sqrt{x^2+y^2}\to0.
\]
所以 \(f\) 在 \((0,0)\) 可微，且全微分为 \(0\)。

再看 \(f_x\) 的连续性。由例11，当 \(xy\ne0\) 时
\[
f_x=2x\sin\frac1{xy}
-\frac{x^2+y^2}{x^2y}\cos\frac1{xy}.
\]
该式在 \((x,y)\to(0,0)\) 时不仅不趋于 \(0\)，甚至沿适当路径会振荡且无界。因此 \(f_x\) 在 \((0,0)\) 不连续。
\answer{C}
\examnote{“可微”不要求偏导在邻域内连续；偏导连续只是可微的充分条件，不是必要条件。}
\end{solutionblock}

\begin{problemblock}
\textbf{例13.} 设
\[
f(x,y)=
\begin{cases}
xy,&|x|\ge |y|,\\
-xy,&|x|<|y|.
\end{cases}
\]
则 \(f''_{xy}(0,0)\) 和 \(f''_{yx}(0,0)\) 依次为 \(\underline{\hspace{2cm}}\)。
\end{problemblock}
\begin{solutionblock}
\analysis{混合偏导在原点要按定义逐层计算，不能用 Clairaut 定理，因为偏导不连续。}
先求 \(f_x(0,y)\)。若 \(y\ne0\)，在 \(x=0\) 附近有 \(|x|<|y|\)，故
\[
f(x,y)=-xy,
\]
于是
\[
f_x(0,y)=-y.
\]
又 \(f_x(0,0)=0\)。因此
\[
f''_{xy}(0,0)
=\lim_{y\to0}\frac{f_x(0,y)-f_x(0,0)}y
=\lim_{y\to0}\frac{-y}{y}=-1.
\]
再求 \(f_y(x,0)\)。若 \(x\ne0\)，在 \(y=0\) 附近有 \(|x|\ge|y|\)，故
\[
f(x,y)=xy,
\]
于是
\[
f_y(x,0)=x.
\]
又 \(f_y(0,0)=0\)。因此
\[
f''_{yx}(0,0)
=\lim_{x\to0}\frac{f_y(x,0)-f_y(0,0)}x
=1.
\]
\answer{\(-1,\ 1\)}
\examnote{二阶混合偏导相等需要连续性条件；本题正是反例。}
\end{solutionblock}

\section{方程化简与多元极值}

\begin{problemblock}
\textbf{例14.} 已知
\[
z(x,y)=xy\,f[\ln(xy)]-(xy)^2,
\]
且 \(f(t)\) 可导，\(z(1,1)=1\)，若
\[
\frac1y\frac{\partial z}{\partial x}
+\frac1x\frac{\partial z}{\partial y}=0,
\]
则（\quad）

A. \(f(1)=\dfrac1e,\ f'(1)=-\dfrac1e\)

B. \(f(1)=e-\dfrac1e,\ f'(1)=e+\dfrac1e\)

C. \(f(1)=\dfrac4e,\ f'(1)=2-\dfrac4e\)

D. \(f(1)=e+\dfrac1e,\ f'(1)=e-\dfrac1e\)
\end{problemblock}
\begin{solutionblock}
\analysis{\(z\) 只依赖于 \(s=xy\)。把偏微分方程化成关于 \(f(t)\) 的一阶微分方程。}
令
\[
s=xy,\qquad t=\ln s.
\]
则
\[
z=s f(\ln s)-s^2.
\]
因为
\[
z_x=z_s y,\qquad z_y=z_s x,
\]
所以
\[
\frac1y z_x+\frac1x z_y=2z_s.
\]
题设给出 \(z_s=0\)。而
\[
z_s=f(t)+f'(t)-2s=f(t)+f'(t)-2e^t.
\]
故
\[
f'(t)+f(t)=2e^t.
\]
又 \(z(1,1)=1\)，此时 \(s=1,t=0\)，
\[
f(0)-1=1,\qquad f(0)=2.
\]
解微分方程：
\[
(e^t f(t))'=2e^{2t}.
\]
积分得
\[
e^t f(t)=e^{2t}+C,\qquad f(t)=e^t+Ce^{-t}.
\]
由 \(f(0)=2\) 得 \(C=1\)，所以
\[
f(t)=e^t+e^{-t},\qquad f'(t)=e^t-e^{-t}.
\]
于是
\[
f(1)=e+\frac1e,\qquad f'(1)=e-\frac1e.
\]
\answer{D}
\examnote{若函数只通过 \(xy\) 进入，常令 \(s=xy\)，再把偏导组合改写为关于 \(s\) 的导数。}
\end{solutionblock}

\begin{problemblock}
\textbf{例15.} 已知函数 \(f(u)\) 在区间 \((0,+\infty)\) 具有二阶导数，记
\[
g(x,y)=f\left(\frac{x}{y}\right).
\]
若 \(g(x,y)\) 满足
\[
x^2g_{xx}+xyg_{xy}+y^2g_{yy}=1,
\]
且
\[
g(x,x)=1,\qquad
\left.\frac{\partial g}{\partial x}\right|_{(x,x)}=\frac2x,
\]
求 \(f(u)\)。
\end{problemblock}
\begin{solutionblock}
\analysis{设 \(u=x/y\)，把 PDE 全部化为 \(u\) 的常微分方程。}
由
\[
g(x,y)=f(u),\qquad u=\frac{x}{y},
\]
可得
\[
g_x=\frac{f'(u)}y,\qquad g_{xx}=\frac{f''(u)}{y^2},
\]
\[
g_{xy}=-\frac{u f''(u)+f'(u)}{y^2},
\]
\[
g_{yy}=\frac{u^2f''(u)+2u f'(u)}{y^2}.
\]
代入 PDE：
\[
x^2g_{xx}+xyg_{xy}+y^2g_{yy}
=u^2f''(u)+u f'(u)=1.
\]
即
\[
u^2f''(u)+u f'(u)=1.
\]
这可写成
\[
u(uf'(u))'=1,
\]
所以
\[
(uf'(u))'=\frac1u.
\]
积分：
\[
uf'(u)=\ln u+C_1,
\]
\[
f(u)=\frac12(\ln u)^2+C_1\ln u+C_2.
\]
由 \(g(x,x)=f(1)=1\)，得
\[
C_2=1.
\]
又
\[
g_x(x,x)=\frac{f'(1)}x=\frac2x,
\]
故
\[
f'(1)=2.
\]
由 \(uf'(u)=\ln u+C_1\) 得 \(C_1=2\)。因此
\[
f(u)=1+2\ln u+\frac12(\ln u)^2.
\]
\answer{\(\displaystyle f(u)=1+2\ln u+\frac12(\ln u)^2\)}
\examnote{齐次型 \(g(x,y)=f(x/y)\) 的偏导组合通常会化为关于 \(u=x/y\) 的 Euler 型方程。}
\end{solutionblock}

\begin{problemblock}
\textbf{例16.} \(z(x,y)\) 具有二阶连续偏导数，设变换
\[
\begin{cases}
u=x-2y,\\
v=x+ay
\end{cases}
\]
可把方程
\[
6z_{xx}+z_{xy}-z_{yy}=0
\]
简化为
\[
z_{uv}=0.
\]
求常数 \(a\)。
\end{problemblock}
\begin{solutionblock}
\analysis{把 \(\partial_x,\partial_y\) 用 \(\partial_u,\partial_v\) 表示，再让 \(z_{vv}\) 项消失，同时变换不能退化。}
由
\[
u=x-2y,\qquad v=x+ay,
\]
得
\[
\partial_x=\partial_u+\partial_v,\qquad
\partial_y=-2\partial_u+a\partial_v.
\]
于是
\[
z_{xx}=z_{uu}+2z_{uv}+z_{vv},
\]
\[
z_{xy}=-2z_{uu}+(a-2)z_{uv}+az_{vv},
\]
\[
z_{yy}=4z_{uu}-4az_{uv}+a^2z_{vv}.
\]
代入：
\[
6z_{xx}+z_{xy}-z_{yy}
=(10+5a)z_{uv}+(6+a-a^2)z_{vv}.
\]
要化为 \(z_{uv}=0\)，需
\[
6+a-a^2=0.
\]
解得
\[
a=3\quad\text{或}\quad a=-2.
\]
但 \(a=-2\) 时 \(v=x-2y=u\)，变换退化，不可取。因此
\[
a=3.
\]
\answer{\(3\)}
\examnote{变量变换化简 PDE 时，除了消掉多余项，还要检查变换的 Jacobian 不能为零。}
\end{solutionblock}

\begin{problemblock}
\textbf{例17.} 设 \(f(x,y)\) 是一阶偏导数连续的正值函数，满足
\[
f_x'(x,y)+f(x,y)=0.
\]
若
\[
f_y'(0,y)=\tan y,\qquad f(0,0)=1,
\]
求 \(f(x,y)\)。
\end{problemblock}
\begin{solutionblock}
\analysis{把 \(y\) 看作参数，先解关于 \(x\) 的一阶线性方程。}
由
\[
f_x+f=0
\]
得
\[
f(x,y)=C(y)e^{-x}.
\]
于是
\[
f_y(0,y)=C'(y)=\tan y.
\]
积分：
\[
C(y)=\int \tan y\,dy=-\ln|\cos y|+C_0.
\]
在 \(y=0\) 附近 \(\cos y>0\)，且
\[
f(0,0)=C(0)=1,
\]
所以 \(C_0=1\)。故
\[
f(x,y)=e^{-x}\bigl(1-\ln(\cos y)\bigr).
\]
\answer{\(\displaystyle f(x,y)=e^{-x}\bigl(1-\ln(\cos y)\bigr)\)}
\examnote{偏微分方程 \(f_x+f=0\) 可先按 \(x\) 解，积分常数变成 \(y\) 的函数。}
\end{solutionblock}

\begin{problemblock}
\textbf{例18.} 设 \(z=z(u,v)\) 具有二阶连续偏导数，且 \(z=z(x-y,x+2y)\) 满足
\[
2z_{xx}+z_{xy}-z_{yy}=2z_x-z_y,
\]
\[
\frac{d[z(0,v)]}{dv}=\frac13z(0,v)+e^{v/3},\qquad z(u,0)=\sin u.
\]
\begin{enumerate}[label=(\arabic*), leftmargin=2.2em]
\item 证明
\[
z_{uv}=\frac13z_u;
\]
\item 求 \(z=z(u,v)\) 的表达式。
\end{enumerate}
\end{problemblock}
\begin{solutionblock}
\analysis{先用 \(u=x-y,\ v=x+2y\) 把 PDE 化为 \(u,v\) 形式；再把 \(z_u\) 看成未知函数解一阶方程。}
(1) 由
\[
u=x-y,\qquad v=x+2y,
\]
得
\[
\partial_x=\partial_u+\partial_v,\qquad
\partial_y=-\partial_u+2\partial_v.
\]
计算
\[
z_{xx}=z_{uu}+2z_{uv}+z_{vv},
\]
\[
z_{xy}=-z_{uu}+z_{uv}+2z_{vv},
\]
\[
z_{yy}=z_{uu}-4z_{uv}+4z_{vv}.
\]
所以
\[
2z_{xx}+z_{xy}-z_{yy}=9z_{uv}.
\]
又
\[
2z_x-z_y=2(z_u+z_v)-(-z_u+2z_v)=3z_u.
\]
故
\[
9z_{uv}=3z_u,
\qquad
z_{uv}=\frac13z_u.
\]

(2) 令
\[
p=z_u.
\]
由 \(p_v=\frac13p\)，得
\[
p(u,v)=A(u)e^{v/3}.
\]
对 \(u\) 积分：
\[
z(u,v)=e^{v/3}B(u)+C(v).
\]
由 \(z(u,0)=\sin u\)，得
\[
B(u)+C(0)=\sin u.
\]
设 \(C(0)=c\)，则
\[
B(u)=\sin u-c.
\]
于是
\[
z(0,v)=-ce^{v/3}+C(v).
\]
题设
\[
\frac d{dv}z(0,v)=\frac13z(0,v)+e^{v/3}
\]
代入后，含 \(c\) 的项抵消，得到
\[
C'(v)=\frac13C(v)+e^{v/3}.
\]
解得
\[
C(v)=(v+c)e^{v/3}.
\]
因此
\[
z(u,v)=e^{v/3}(\sin u-c)+(v+c)e^{v/3}
=e^{v/3}(\sin u+v).
\]
\answer{(1) 已证；(2) \(\displaystyle z(u,v)=e^{v/3}(\sin u+v)\)。}
\examnote{先把 PDE 化成 \(z_u\) 的一阶方程，再用边界条件确定积分函数。}
\end{solutionblock}

\begin{problemblock}
\textbf{例19.} 设
\[
f(x,y)=(y-x^2)(y-2x^2),
\]
\(k\) 为任意常数，则（\quad）

A. \(f(x,kx)\) 在 \(x=0\) 处取极小值，点 \((0,0)\) 是 \(f(x,y)\) 的极小值点

B. \(f(x,kx)\) 在 \(x=0\) 处取极小值，点 \((0,0)\) 不是 \(f(x,y)\) 的极小值点

C. \(f(x,kx)\) 在 \(x=0\) 处不取极小值，点 \((0,0)\) 是 \(f(x,y)\) 的极小值点

D. \(f(x,kx)\) 在 \(x=0\) 处不取极小值，点 \((0,0)\) 不是 \(f(x,y)\) 的极小值点
\end{problemblock}
\begin{solutionblock}
\analysis{沿任意直线看是极小，不代表二元函数有极小；还要检查曲线路径。}
沿直线 \(y=kx\)，
\[
f(x,kx)=(kx-x^2)(kx-2x^2)=x^2(k-x)(k-2x).
\]
若 \(k\ne0\)，则当 \(x\) 充分接近 \(0\) 时，
\[
(k-x)(k-2x)>0,
\]
故 \(f(x,kx)\ge0=f(0,0)\)，在 \(x=0\) 处取极小值。若 \(k=0\)，则
\[
f(x,0)=2x^4\ge0,
\]
也在 \(x=0\) 处取极小值。

但对二元函数，取曲线路径
\[
y=\frac32x^2,
\]
则
\[
f\left(x,\frac32x^2\right)
=\left(\frac12x^2\right)\left(-\frac12x^2\right)
=-\frac14x^4<0\quad(x\ne0).
\]
而 \(f(0,0)=0\)，故 \((0,0)\) 不是 \(f(x,y)\) 的极小值点。
\answer{B}
\examnote{多元极值不能只看所有直线路径；曲线路径可能揭示相反符号。}
\end{solutionblock}

\begin{problemblock}
\textbf{例20.} 设 \(f(x,y)\) 与 \(\varphi(x,y)\) 均为可微函数，且
\[
\varphi_y'(x,y)\ne0.
\]
已知 \((x_0,y_0)\) 是 \(f(x,y)\) 在约束条件 \(\varphi(x,y)=0\) 下的一个极值点，下列选项正确的是（\quad）

A. 若 \(f_y'(x_0,y_0)=0\)，则 \(f_x'(x_0,y_0)=0\)

B. 若 \(f_x'(x_0,y_0)=0\)，则 \(f_y'(x_0,y_0)\ne0\)

C. 若 \(f_x'(x_0,y_0)\ne0\)，则 \(f_y'(x_0,y_0)=0\)

D. 若 \(f_x'(x_0,y_0)\ne0\)，则 \(f_y'(x_0,y_0)\ne0\)
\end{problemblock}
\begin{solutionblock}
\analysis{由 \(\varphi_y\ne0\)，约束曲线可写成 \(y=y(x)\)。约束极值要求一元复合函数导数为零。}
由隐函数定理，\(\varphi(x,y)=0\) 可在点附近写为
\[
y=y(x),
\qquad
y'(x_0)=-\frac{\varphi_x'(x_0,y_0)}{\varphi_y'(x_0,y_0)}.
\]
令
\[
F(x)=f(x,y(x)).
\]
约束极值给出
\[
F'(x_0)=f_x'(x_0,y_0)+f_y'(x_0,y_0)y'(x_0)=0.
\]
等价地，也可用 Lagrange 乘数：
\[
\nabla f(x_0,y_0)=\lambda\nabla\varphi(x_0,y_0).
\]
由于 \(\varphi_y'(x_0,y_0)\ne0\)，若
\[
f_y'(x_0,y_0)=0,
\]
则 \(\lambda=0\)，从而
\[
f_x'(x_0,y_0)=0.
\]
所以 A 正确。其逆否命题为
\[
f_x'(x_0,y_0)\ne0\Rightarrow f_y'(x_0,y_0)\ne0,
\]
即 D 也正确。B、C 与此相反，错误。
\answer{A、D}
\examnote{原题选项 A 与 D 是互为逆否的同一逻辑结论；若按单选题排版，这是题面重复正确项。}
\end{solutionblock}

\begin{problemblock}
\textbf{例21.} 函数
\[
f(x,y)=x^2+y^2
\]
在约束条件
\[
(x-1)^3=y^2
\]
下（\quad）

A. 有最大值，无最小值\qquad
B. 无最大值，无最小值

C. 有最大值，有最小值\qquad
D. 无最大值，有最小值
\end{problemblock}
\begin{solutionblock}
\analysis{约束曲线是半立方抛物线，可参数化为 \(x=1+t^2,\ y=t^3\)。}
由
\[
(x-1)^3=y^2
\]
可设
\[
x=1+t^2,\qquad y=t^3,\qquad t\in\mathbb{R}.
\]
于是
\[
f=x^2+y^2=(1+t^2)^2+t^6
=1+2t^2+t^4+t^6.
\]
显然
\[
f\ge1,
\]
且当 \(t=0\)，即点 \((1,0)\) 时取到最小值 \(1\)。

当 \(|t|\to+\infty\) 时，
\[
f\to+\infty,
\]
故无最大值。
\answer{D}
\examnote{非闭有界约束集上，连续函数不一定有最大值；先参数化判断是否无界。}
\end{solutionblock}

\begin{problemblock}
\textbf{例22.} 已知曲线
\[
C:\begin{cases}
x^2+y^2-2z-12=0,\\
x+y+z=3,
\end{cases}
\]
求曲线 \(C\) 上距离原点最远点和最近点。
\end{problemblock}
\begin{solutionblock}
\analysis{先用平面方程消去 \(z\)，把空间曲线化为平面内圆，再优化距离平方。}
由
\[
z=3-x-y
\]
代入
\[
x^2+y^2-2z-12=0
\]
得
\[
x^2+y^2+2x+2y-18=0,
\]
即
\[
(x+1)^2+(y+1)^2=20.
\]
令
\[
a=x+1,\qquad b=y+1.
\]
则
\[
a^2+b^2=20,\qquad z=5-a-b.
\]
距离平方为
\[
\begin{aligned}
R^2&=x^2+y^2+z^2\\
&=(a-1)^2+(b-1)^2+(5-a-b)^2.
\end{aligned}
\]
化简得
\[
R^2=(a+b)^2-12(a+b)+47.
\]
设
\[
s=a+b.
\]
在约束 \(a^2+b^2=20\) 下，
\[
-2\sqrt{10}\le s\le2\sqrt{10}.
\]
于是
\[
R^2=(s-6)^2+11.
\]
最近时 \(s=6\)，此时
\[
a+b=6,\qquad a^2+b^2=20.
\]
解得
\[
(a,b)=(2,4)\quad\text{或}\quad(4,2).
\]
对应
\[
(x,y,z)=(1,3,-1),\ (3,1,-1).
\]
最远时 \(s=-2\sqrt{10}\)，此时 \(a=b=-\sqrt{10}\)，对应
\[
x=y=-1-\sqrt{10},\qquad z=5+2\sqrt{10}.
\]
\answer{最近点为 \((1,3,-1)\)、\((3,1,-1)\)；最远点为 \((-1-\sqrt{10},-1-\sqrt{10},5+2\sqrt{10})\)。}
\examnote{空间曲线距离问题通常先消元，优化距离平方比优化距离本身更简单。}
\end{solutionblock}

\begin{problemblock}
\textbf{例23.} 求函数
\[
f(x,y)=3(x^2+y^2)-x^3
\]
在有界闭区域
\[
D=\{(x,y)\mid x^2+y^2\le16\}
\]
上的最大值和最小值。
\end{problemblock}
\begin{solutionblock}
\analysis{闭区域上连续函数必有最大最小值。先查内部驻点，再查边界圆。}
内部驻点由
\[
f_x=6x-3x^2=3x(2-x),\qquad f_y=6y
\]
给出。故内部驻点为
\[
(0,0),\qquad (2,0).
\]
函数值分别为
\[
f(0,0)=0,\qquad f(2,0)=3\cdot4-8=4.
\]
边界为
\[
x^2+y^2=16.
\]
在边界上
\[
f=3\cdot16-x^3=48-x^3,\qquad -4\le x\le4.
\]
因此边界最大值在 \(x=-4\) 处取得：
\[
f_{\max}=48-(-4)^3=112,
\]
对应点
\[
(-4,0).
\]
边界最小值在 \(x=4\) 处取得：
\[
f_{\min}=48-4^3=-16,
\]
对应点
\[
(4,0).
\]
与内部驻点值比较后，整体最大值为 \(112\)，最小值为 \(-16\)。
\answer{最大值 \(112\)，在 \((-4,0)\) 处取得；最小值 \(-16\)，在 \((4,0)\) 处取得。}
\examnote{闭区域极值题固定流程：内部驻点、边界约束、最后统一比较函数值。}
\end{solutionblock}
